\chapter*{Введение}                         % Заголовок
\addcontentsline{toc}{chapter}{Введение}    % Добавляем его в оглавление

\newcommand{\actuality}{}
\newcommand{\progress}{\textbf{\progressTXT}}
\newcommand{\aim}{{\textbf\aimTXT}}
\newcommand{\tasks}{\textbf{\tasksTXT}}
\newcommand{\object}{{\textbf\objectTXT}}
\newcommand{\subject}{{\textbf\subjectTXT}}
\newcommand{\novelty}{\textbf{\noveltyTXT}}
\newcommand{\tinfluence}{\textbf{\tinfluenceTXT}}
\newcommand{\pinfluence}{\textbf{\pinfluenceTXT}}
\newcommand{\methods}{\textbf{\methodsTXT}}
\newcommand{\defpositions}{\textbf{\defpositionsTXT}}
\newcommand{\passport}{{\textbf\passportTXT}}
\newcommand{\reliability}{\textbf{\reliabilityTXT}}
\newcommand{\probation}{\textbf{\probationTXT}}
\newcommand{\contribution}{\textbf{\contributionTXT}}
\newcommand{\publications}{\textbf{\publicationsTXT}}

\providecommand{\extcite}[1]{\ifsynopsis\else~\cite{#1}\fi}% внешние ссылки только в диссертации
{\actuality} В статистической теории обучения обоснована принципиальная связь между сложностью семейства моделей, объемом обучающей выборки и способностью к обобщению\extcite{vapnik1995nature}. Для надежного выбора модели требуется достаточное количество данных. Подавляющее большинство современных моделей машинного обучения являются параметрическими: они задаются конечным набором параметров, определяемых минимизацией эмпирической функции потерь на выборке. Поскольку функция потерь вычисляется по ограниченным данным и зависит от параметров модели, возникает фундаментальный вопрос о взаимосвязи объема выборки, структуры модели и свойств множества оптимальных параметров, а также о \textbf{минимальном объеме данных}, необходимом для успешного обучения\extcite{kaplan2020scaling,hoffmann2022}.

Эмпирические \textbf{законы масштабирования} показывают, что качество моделей закономерно зависит от размера данных и модели: в языковых моделях установлены степенные зависимости потерь от объема данных и числа параметров\extcite{kaplan2020scaling}, а для вычислительно-оптимального обучения предложено согласованное масштабирование размера модели и объема данных\extcite{hoffmann2022}. Эти результаты опираются на масштабные эксперименты и не имеют пока единого теоретического обоснования.

Естественным объектом для такого анализа является \textbf{ландшафт функции потерь} в пространстве параметров. Его визуализация и анализ структуры минимумов\extcite{li2018visualizing}, изучение локальной геометрии и связей с обобщающей способностью\extcite{fort2019large}, а также доказательство связности множества решений\extcite{draxler2018essentially} показывают, что геометрия поверхности функции потерь определяет как выбор решения методами градиентной оптимизации\extcite{garipov2018loss,chaudhari2019}, так и устойчивость к возмущениям данных. Спектральные свойства \textbf{матрицы Гессе} в окрестности минимума характеризуют кривизну оптимизационной поверхности и используются при анализе сходимости и обобщения\extcite{papyan2019}. Вместе с тем остается открытой задача систематического описания того, как ландшафт функции потерь стабилизируется при увеличении объема выборки и как связать эту стабилизацию с законами масштабирования и достаточным размером данных.

\ifsynopsis\textit{Таким образом, актуальной задачей является разработка теоретических методов, связывающих архитектуру параметрической модели, объем обучающей выборки и локальную геометрию множества оптимальных параметров, задаваемую ландшафтом функции потерь.}\else{Таким образом, актуальной задачей является разработка теоретических методов, связывающих архитектуру параметрической модели, объем обучающей выборки и локальную геометрию множества оптимальных параметров, задаваемую ландшафтом функции потерь.}\fi

\ifsynopsis
{\progress} Задача определения достаточного размера выборки традиционно решается методами математической статистики, где размер выборки выбирается так, чтобы критерий проверки гипотез имел заданную мощность при фиксированной вероятности ошибки первого рода, методами байесовского анализа, основанными на апостериорной неопределенности и близости апостериорных распределений, и методами статистической теории обучения, связывающими объем выборки со сложностью класса моделей. Для современных нейронных сетей эти подходы либо требуют явных гипотез о распределении параметров, либо дают сильно завышенные оценки. Ландшафт функции потерь и спектр матрицы Гессе активно исследуются в связи с оптимизацией и обобщающей способностью нейронных сетей, однако связь локальной геометрии оптимизационной поверхности с достаточным размером выборки и законами масштабирования остается открытым вопросом.
\else
% Степень разработанности темы (только в диссертации); перенесено из бывшей главы 1
{\progress} Задача определения достаточного размера обучающей выборки и законы масштабирования изучаются в нескольких областях: в математической статистике, в байесовском анализе, в статистической теории обучения и в исследованиях ландшафта функции потерь нейронных сетей.

Особый интерес представляет вопрос о законах масштабирования параметрических моделей: по какому правилу следует совместно выбирать архитектуру модели и объем данных для ее обучения. В классических постановках задача определения достаточного размера выборки обычно рассматривается для линейных и обобщенно-линейных моделей, где модель задается конечномерным вектором параметров, а достаточность выборки формулируется через статистические критерии различимости гипотез о параметрах или через устойчивость вероятностной оценки модели~\cite{self1988,self1992,shieh2000,shieh2005,demidenko2007,motrenko2014}. Для обобщенно-линейных моделей используются, в частности, нормальная, логистическая и пуассоновская регрессии, объединяемые общей экспоненциальной формой распределения отклика и функцией связи между математическим ожиданием и линейным предиктором.

Классические статистические методы оценки размера выборки опираются на задание нулевой и альтернативной гипотез, допустимой вероятности ошибки первого рода и требуемой мощности теста. В наиболее общем виде достаточный размер выборки $m^*$ определяется как минимальное значение $m$, при котором статистика критерия $T_m$ обеспечивает требуемую мощность:
\[
    m^*:\quad \mathbb P(T_m > c_\alpha | H_1)\ge 1-\beta,
\]
где $c_\alpha$ "--- критическое значение уровня значимости $\alpha$, а $(1-\beta)$ "--- требуемая мощность. В этой постановке используются критерий отношения правдоподобия, статистика Вальда, метод множителей Лагранжа и их модификации~\cite{self1988,self1992,shieh2000,shieh2005,demidenko2007}. Эти методы дают теоретически обоснованные оценки, но требуют явных предположений о распределении данных, структуре модели и величине эффекта при альтернативе. На практике такие предположения часто недоступны или труднопроверяемы, особенно для многопараметрических нелинейных моделей.

Наряду со статистическими существуют байесовские и информационно-теоретические подходы. В байесовской постановке достаточность выборки связывается со стабилизацией апостериорного распределения параметров, уменьшением апостериорной неопределенности или малостью расстояния между распределениями, полученными по соседним подвыборкам. Типичными критериями являются средняя апостериорная дисперсия, среднее покрытие, средняя длина доверительной области и дивергенция Кульбака--Лейблера между апостериорными распределениями~\cite{lindley1997,rubin1998,wang2002,motrenko2014}. Информационно-теоретические подходы близки к этой линии и также используют энтропийные и дивергентные характеристики для описания связи между моделью и данными. Однако для моделей большой размерности такие методы часто оказываются вычислительно сложными и требуют дополнительных модельных предположений.

В статистической теории обучения размер выборки связывается со сложностью класса моделей через размерность Вапника"--~Червоненкиса (англ. Vapnik"--~Chervonenkis dimension, VC) и родственные меры емкости. Классические результаты этого типа лежат в основе анализа вероятно приблизительно корректного обучения (англ. probably approximately correct, PAC) и дают оценки числа объектов, достаточного для гарантированного обобщения. Однако для нейронных сетей VC-размерность, как правило, очень велика, вследствие чего получаемые оценки оказываются грубыми и существенно завышенными. Более современные меры, зависящие от выборки, такие как радемахеровские сложности, позволяют получать более тонкие оценки обобщающей способности. В русскоязычной литературе существенный вклад в развитие этой линии внес К.~В.~Воронцов, развивший комбинаторную теорию переобучения, методы получения точных оценок вероятности переобучения и оценки обобщающей способности, зависящие от выборки, в том числе в связи с радемахеровской сложностью и комбинаторными оценками надежности обучения по прецедентам~\cite{vorontsov2008pria,vorontsov2010dsc,vorontsov2012overfitting}. Для глубоких сетей были также предложены оценки через спектральную норму и отступ (англ. margin bounds), а также их уточнения, учитывающие произведения операторных норм весовых матриц~\cite{bartlett2017spectral}. Тем не менее даже такие оценки далеки от объяснения высокой обобщающей способности сильно перепараметризованных сетей, наблюдаемой на практике. Экспериментально показано, что современные глубокие сети способны точно аппроксимировать даже случайную разметку, что плохо согласуется с наивными комбинаторными оценками сложности~\cite{zhang2017rethinking}.

Классические подходы плохо переносятся на современные нейронные сети по нескольким причинам. Во-первых, комбинаторные и асимптотические меры сложности дают нереалистично большие или трудно интерпретируемые оценки необходимого объема данных. Во-вторых, явные статистические гипотезы о параметрах нейросетей обычно неизвестны. В-третьих, многократное полное переобучение модели на подвыборках разного размера для построения кривых обучения вычислительно дорого для больших сетей.

В то же время эмпирические исследования последних лет показывают существование устойчивых законов масштабирования: качество моделей систематически улучшается с ростом числа параметров, объема данных и вычислительного бюджета. В работе~\cite{kaplan2020scaling} показано, что для языковых моделей значение функции потерь хорошо аппроксимируется степенными законами по размеру модели, размеру выборки и объему вычислений. В работе~\cite{hoffmann2022} предложена постановка, оптимальная по вычислительному бюджету (англ. compute-optimal), согласно которой при фиксированном бюджете размер модели и число обучающих элементов последовательности (англ. tokens) должны масштабироваться согласованно; при этом многие крупные языковые модели оказались недообученными по данным. Эти результаты оказали значительное влияние на современное понимание компромисса между размером модели и объемом обучающей выборки, однако единого теоретического объяснения таких закономерностей для широкого класса параметрических моделей до сих пор нет.

Таким образом, возникает необходимость в развитии новых подходов к описанию достаточного размера выборки и законов масштабирования, опирающихся на свойства, доступные для анализа во всех параметрических моделях и инвариантные к выбору конкретного минимума. В настоящей диссертации в качестве такого объекта предлагается использовать локальную геометрию задачи минимизации эмпирической функции потерь, то есть свойства ландшафта функции потерь в окрестности точки минимума. Метод оценки масштабируемости, который использует эту локальную геометрию, а не глобальную картину поверхности, называется ландшафтным.

В задачах машинного обучения оптимум эмпирической функции потерь часто не единствен. В линейной регрессии при невырожденной матрице признаков существует единственный глобальный минимум, но при мультиколлинеарности или при $m < n$ задача вырождена: множество решений образует аффинное подпространство~\cite{seber2003,strang2006,vorontsov2010dsc}. В нейронных сетях функция потерь, как правило, невыпукла: существуют множественные локальные минимумы, седловые точки и плато~\cite{dauphin2014,choromanska2015,goodfellow2015qualitatively,li2018visualizing}. Более того, из-за симметрий (перестановка нейронов, знаковые преобразования) многие минимумы эквивалентны с точки зрения предсказательной способности~\cite{draxler2018essentially,entezari2022role,ainsworth2023git}. В таких условиях корректнее говорить не о единственной точке минимума, а о множестве оптимальных параметров модели.

Оптимизационной поверхностью или ландшафтом функции потерь (англ. loss landscape) называется график эмпирической функции потерь $\mathcal{L}_m(\mathbf{w})$, вычисленной по выборке размера $m$, как функции от параметров $\mathbf{w} \in \mathbb{R}^N$. Анализ этой поверхности для нейронных сетей является ключевой задачей теории оптимизации, что связано с итерационной природой процесса обучения "--- алгоритмы на основе стохастического градиентного спуска (англ. stochastic gradient descent, SGD) дают приближенное решение задачи минимизации эмпирической функции потерь, опираясь на глобальные и локальные свойства.

Исследование ландшафта функции потерь позволило выявить ряд важных характеристик, определивших развитие современных методов численной оптимизации~\cite{kingma2015adam,chaudhari2019,foret2021sam} "--- наличие множественных минимумов, седловых точек, многообразий меньшей размерности, содержащих множества оптимальных параметров~\cite{draxler2018essentially,garipov2018loss}, и др. В теории обобщения нейронных сетей плоские минимумы (с малой кривизной) часто связывают с лучшей обобщающей способностью~\cite{chaudhari2019,keskar2017largebatch,foret2021sam}, а узкие и глубокие "--- наоборот. Форма поверхности при этом определяется как самой архитектурой модели, так и заданной обучающей выборкой. Глубина и ширина нейронной сети давно перестали быть определяющими факторами для достижения высокого качества~\cite{he2016resnet}, из-за чего возникают новые архитектурные решения, позволяющие снизить сложность модели при неухудшении ее обобщающих свойств~\cite{ioffe2015batchnorm,ba2016layernorm,ulyanov2016instancenorm,wu2018groupnorm,srivastava2014dropout,henry2020qknorm}.

Ключевой характеристикой для анализа оптимизационной поверхности является матрица Гессе функции потерь $\mathbf{H}^{(m)}(\mathbf{w}) = \nabla^2_{\mathbf{w}} \mathcal{L}_{m}(\mathbf{w})$, которая содержит информацию о поведении ландшафта во втором порядке. Спектр матрицы Гессе характеризует нелинейность поверхности в окрестности минимума: большие собственные значения соответствуют  <<острым>> направлениям, малые "--- <<плоским>>~\cite{hochreiter1997flat,dauphin2014,keskar2017largebatch,li2018visualizing}. Более того, для типичных нейросетей наблюдается характерная структура спектра: множество малых собственных значений и небольшое число выбросов~\cite{papyan2019,ghorbani2019investigation,sankar2021deeper,singh2021analytic}. Множество работ посвящено еще более детальному анализу матрицы Гессе для различных архитектур моделей: выводу его явных выражений для полносвязных сетей~\cite{bishop1992hessian,pearlmutter1994,martens2012curvprop} и даже механизмов внимания в трансформерных сетях~\cite{ormaniec2024transformerhessian}.

Матрицы Гессе используются в методе Ньютона и в квазиньютоновских методах~\cite{nocedal2006numerical,dennis1983}. Наиболее известен среди последних алгоритм Бройдена"--~Флетчера"--~Гольдфарба"--~Шанно (англ. Broyden"--~Fletcher"--~Goldfarb"--~Shanno, BFGS)~\cite{broyden1970,fletcher1970}. Матрица Гессе и произведения Гессе на вектор применяются не только для построения методов второго порядка, но и для анализа локальной кривизны ландшафта, оценки <<остроты>> и <<плоскости>> минимумов, спектрального исследования динамики обучения, приближенного вычисления доверительных характеристик параметров, прореживания сети (англ. pruning) и малоранговых приближений (англ. low-rank), а также для построения гаусс"--~ньютоновских методов и методов без явной матрицы Гессе (англ. hessian-free) в больших нейросетевых моделях~\cite{martens2010hf,martens2012curvprop,papyan2019,zhang2024whytransformers}.

Несмотря на интерес научного сообщества к исследованию матрицы Гессе и ее применений в обучении нейронных сетей, вопрос о связи характеристик второго порядка оптимизационной поверхности с законами масштабирования остается открытым.

\fi

{\aim} данной работы является разработка ландшафтного метода оценки масштабируемости параметрических моделей машинного обучения, основанного на локальной геометрии ландшафта функции потерь в окрестности точки минимума и на стабилизации распределения оптимальных параметров. Для достижения цели были поставлены и решены следующие {\tasks}:

\begin{enumerate}[beginpenalty=10000]
    \item Исследовать спектральные свойства матриц Гессе нейронных сетей и получить архитектурно-зависимые верхние оценки спектральных норм для полносвязных и сверточных архитектур.
    \item Разработать критерии стабилизации оптимизационной поверхности при увеличении объема обучающей выборки.
    \item Оценить масштабируемость параметрических моделей машинного обучения на основе локальной геометрии ландшафта функции потерь, связав ее с архитектурой модели и объемом обучающей выборки.
    \item Построить критерии достаточного размера выборки для параметрических моделей на основе стабилизации ландшафта функции потерь.
    \item Исследовать вероятностные линейные модели как частный случай параметрических моделей и построить для них критерии достаточного размера выборки на основе сходимости распределения оптимальных параметров.
\end{enumerate}

{\object} настоящей работы являются параметрические модели машинного обучения и возникающие при их обучении оптимизационные задачи на конечных обучающих выборках.

{\subject} данной работы являются локальные свойства оптимизационной поверхности функции потерь параметрических моделей машинного обучения, спектральные свойства матрицы Гессе, стабилизация множества оптимальных параметров при увеличении объема обучающей выборки и критерии достаточного размера выборки, основанные на этой стабилизации.

{\novelty}
\begin{enumerate}[beginpenalty=10000]
    \item Получены верхние оценки спектральных норм гаусс-ньютоновской компоненты матрицы Гессе для полносвязных и сверточных архитектур нейронных сетей на основе ее блочной факторизации в классе матрично-представимых моделей.
    \item Предложены критерии стабилизации оптимизационной поверхности. Порядок их убывания следует из усреднения эмпирического риска по объектам выборки; локальная квадратичная аппроксимация и спектр матрицы Гессе оценивают константу в этой оценке.
    \item Предложена схема законов масштабирования, в которой ошибка оценивания задается отношением шума градиентов к кривизне. Для линейной модели с неполным набором признаков схема выполнена точно.
    \item Получены критерии достаточного размера выборки для параметрических моделей на основе стабилизации ландшафта функции потерь.
    \item Для вероятностных линейных моделей построены критерии достаточного размера выборки, основанные на сходимости апостериорного распределения параметров и стабилизации вероятностных характеристик решения.
\end{enumerate}

{\tinfluence} полученных результатов состоит в том, что предложено обоснование связи между объемом обучающей выборки и стабилизацией оптимизационной поверхности (для линейных моделей "--- через правдоподобие и апостериорные распределения, для нейронных сетей "--- через спектр матрицы Гессе и сходимость ландшафта функции потерь), дополняющее известные эмпирические законы масштабирования строгими оценками и единой постановкой задачи.

{\pinfluence} настоящей работы заключается в возможности применения разработанных методов для обоснованного определения достаточного объема данных при планировании экспериментов и сборе выборок для линейных и нейросетевых моделей, а также для формулировки критериев остановки сбора данных и интерпретации достаточности выборки в терминах геометрии поверхности функции потерь.

{\methods} Для достижения поставленных
в диссертационной работе целей используются:
\begin{enumerate}[beginpenalty=10000]
    \item Методы линейной алгебры, математического и функционального анализа: матричное дифференцирование, анализ функций многих переменных, спектральный анализ матриц (в том числе матрицы Гессе).
    \item Методы теории вероятностей и математической статистики: байесовский вывод, анализ сходимости последовательностей случайных величин и распределений, элементы теории оптимизации.
    \item Методы вычислительного эксперимента и анализа результатов: бутстрап (англ. bootstrap), метод Монте"--~Карло, воспроизводимая постановка экспериментов.
\end{enumerate}

{\defpositions}
\begin{enumerate}[beginpenalty=10000]
    \item Теоремы о верхних оценках спектральной нормы гаусс-ньютоновской компоненты матрицы Гессе для полносвязных и сверточных архитектур нейронных сетей.
    \item Критерии стабилизации ландшафта функции потерь: порядок $\mathcal O(k^{-1})$ и $\mathcal O(k^{-2})$ следует из усреднения по объектам, а оценки констант получены по локальной квадратичной аппроксимации в одной точке, в полном пространстве параметров и в подпространстве главных собственных направлений матрицы Гессе.
    \item Схема законов масштабирования, в которой ошибка оценивания задается отношением шума градиентов к кривизне, и ее точный вид для линейной модели с неполным набором признаков.
    \item Критерии достаточного размера выборки для параметрических моделей, основанные на стабилизации ландшафта функции потерь.
    \item Критерии достаточного размера выборки для вероятностных линейных моделей, основанные на стабилизации распределения оптимальных параметров при увеличении объема данных.
\end{enumerate}

{\passport} Диссертационная работа соответствует паспорту научной специальности 1.2.1 "--- <<Искусственный интеллект и машинное обучение>>, а именно в наибольшей степени следующим направлениям исследований: п.~1 <<Естественно-научные основы и методы искусственного интеллекта>>; п.~2 <<Исследования в области оценки качества и эффективности алгоритмических решений для систем искусственного интеллекта и машинного обучения>>; п.~12 <<Исследования в области надежности, устойчивости и переобучения систем класса ИИ>>; п.~15 <<Математические исследования в области статистики, логики, алгебры, топологии, анализа функций и других областях, ориентированные на решение задач искусственного интеллекта и машинного обучения>>; п.~16 <<Исследования в области специальных методов оптимизации, проблем сложности и снижения размерности>>; п.~17 <<Исследования в области многослойных алгоритмических конструкций, в том числе многослойных нейросетей>>.

{\reliability} полученных результатов обеспечивается корректностью
применения апробированного в научной практике математического аппарата теории вероятностей, математической статистики и других смежных теоретических областей, а также экспериментальной проверкой разработанных методов на большом количестве модельных и практических задач машинного обучения. Полученные теоретические результаты обосновываются математически строгими доказательствами, а для проведенных вычислительных экспериментов даются детальные описания, обеспечивающие их воспроизводимость. Результаты опубликованы в рецензируемых научных изданиях и в трудах рецензируемых российских и международных конференций по машинному обучению и искусственному интеллекту.

{\probation} Основные результаты работы докладывались и обсуждались на следующих научных конференциях:
\begin{enumerate}
    \item 16-я Международная конференция
    <<Интеллектуализация обработки информации>> (ИОИ-2026),
    Витебск, 28~сентября~-- 2~октября 2026~г.
    \item Международная конференция <<Иванниковские чтения>>,
    Этномир, 28--30~мая 2026~г.
    \item 68-я Всероссийская научная конференция МФТИ,
    Москва, 30~марта~-- 4~апреля 2026~г.
    \item Открытая конференция ИСП РАН, Москва, 9--10 декабря 2025~г.
    \item 22-я Всероссийская конференция с международным участием
    <<Математические методы распознавания образов (ММРО-2025)>>,
    Муром, 22--26 сентября 2025~г.
    \item 67-я Всероссийская научная конференция МФТИ,
    Москва, 31 марта -- 5 апреля 2025~г.
    \item Открытая конференция ИСП РАН, Москва, 11--12 декабря 2024~г.
    \item 66-я Всероссийская научная конференция МФТИ,
    Москва, 1--6 апреля 2024~г.
\end{enumerate}


{\contribution} Все приведенные результаты, кроме отдельно оговоренных случаев, получены диссертантом лично при научном руководстве к.ф.-м.н.~А.\,В.~Грабового.

\ifnumequal{\value{bibliosel}}{0}
{%%% Встроенная реализация с загрузкой файла через движок bibtex8. (При желании, внутри можно использовать обычные ссылки, наподобие `\cite{vakbib1,vakbib2}`).
    {\publications} Основные результаты по теме диссертации изложены
    в~XX~печатных изданиях,
    X из которых изданы в журналах, рекомендованных ВАК,
    X "--- в тезисах докладов.
}%
{%%% Реализация пакетом biblatex через движок biber
    \begin{refsection}[bl-author, bl-registered]
        % Это refsection=1.
        % Процитированные здесь работы:
        %  * подсчитываются, для автоматического составления фразы "Основные результаты ..."
        %  * попадают в авторскую библиографию, при usefootcite==0 и стиле `\insertbiblioauthor` или `\insertbiblioauthorgrouped`
        %  * нумеруются там в зависимости от порядка команд `\printbibliography` в этом разделе.
        %  * при использовании `\insertbiblioauthorgrouped`, порядок команд `\printbibliography` в нем должен быть тем же (см. biblio/biblatex.tex)
        %
        % Невидимый библиографический список для подсчета количества публикаций:
        \phantom{\printbibliography[heading=nobibheading, section=1, env=countauthorvak,          keyword=biblioauthorvak]%
        \printbibliography[heading=nobibheading, section=1, env=countauthorwos,          keyword=biblioauthorwos]%
        \printbibliography[heading=nobibheading, section=1, env=countauthorscopus,       keyword=biblioauthorscopus]%
        \printbibliography[heading=nobibheading, section=1, env=countauthorconf,         keyword=biblioauthorconf]%
        \printbibliography[heading=nobibheading, section=1, env=countauthorother,        keyword=biblioauthorother]%
        \printbibliography[heading=nobibheading, section=1, env=countauthor,             keyword=biblioauthor]%
        \printbibliography[heading=nobibheading, section=1, env=countauthorvakscopuswos, filter=vakscopuswos]%
        \printbibliography[heading=nobibheading, section=1, env=countauthorscopuswos,    filter=scopuswos]}%
        %
        \nocite{*}%
        %
        % {\publications} Основные результаты по теме диссертации изложены в~\arabic{citeauthor}~печатных изданиях~\cite{},
        % \arabic{citeauthorvak} из которых изданы в журналах, рекомендованных ВАК%
        % \ifnum \value{citeauthorscopuswos}>0%
        %     , \arabic{citeauthorscopuswos} "--- в~периодических научных журналах, индексируемых Web of~Science и Scopus%
        % \fi%
        % \ifnum \value{citeauthorconf}>0%
        %     , \arabic{citeauthorconf} "--- в~тезисах докладов.
        % \else%
        %     .
        % \fi%
        % \ifnum \value{citeregistered}=1%
        %     \ifnum \value{citeauthorpatent}=1%
        %         Зарегистрирован \arabic{citeauthorpatent} патент.
        %     \fi%
        %     \ifnum \value{citeauthorprogram}=1%
        %         Зарегистрирована \arabic{citeauthorprogram} программа для ЭВМ.
        %     \fi%
        % \fi%
        % \ifnum \value{citeregistered}>1%
        %     Зарегистрированы\ %
        %     \ifnum \value{citeauthorpatent}>0%
        %     \formbytotal{citeauthorpatent}{патент}{}{а}{}%
        %     \ifnum \value{citeauthorprogram}=0 . \else \ и~\fi%
        %     \fi%
        %     \ifnum \value{citeauthorprogram}>0%
        %     \formbytotal{citeauthorprogram}{программ}{а}{ы}{} для ЭВМ.
        %     \fi%
        % \fi%
        {\publications} \ifsynopsis\textit{Список публикаций приведен в конце автореферата. }\else\textit{Работы автора приведены в общем списке литературы. }\fi Основные результаты по теме диссертации изложены в~\arabic{citeauthor}~научных работах, из которых~\arabic{citeauthorscopuswos} "--- в~рецензируемых периодических научных журналах, а остальные~\arabic{citeauthorconf} "--- в~сборниках трудов конференций и тезисах докладов.
        % К публикациям, в которых излагаются основные научные результаты диссертации на соискание ученой
        % степени, в рецензируемых изданиях приравниваются патенты на изобретения, патенты (свидетельства) на
        % полезную модель, патенты на промышленный образец, патенты на селекционные достижения, свидетельства
        % на программу для электронных вычислительных машин, базу данных, топологию интегральных микросхем,
        % зарегистрированные в установленном порядке.(в ред. Постановления Правительства РФ от 21.04.2016 N 335)
    \end{refsection}%
    \begin{refsection}[bl-author, bl-registered]
        % Это refsection=2.
        % Процитированные здесь работы:
        %  * попадают в авторскую библиографию, при usefootcite==0 и стиле `\insertbiblioauthorimportant`.
        %  * ни на что не влияют в противном случае
        \nocite{kiselev2026ispras}
        \nocite{kiselev2025ssdpdp}
        \nocite{kiselev2025ssdlb}
        \nocite{kiselev2024unraveling}
        \nocite{kiselev2026conf68mipt}
        \nocite{kiselev2025mmro}
        \nocite{kiselev2025conf67mipt}
        \nocite{meshkov2024convnets}
        \nocite{kiselev2024conf66mipt}
    \end{refsection}%
        %
        % Все, что вне этих двух refsection, это refsection=0,
        %  * для диссертации - это нормальные ссылки, попадающие в обычную библиографию
        %  * для автореферата:
        %     * при usefootcite==0, ссылка корректно сработает только для источника из `external.bib`. Для своих работ "--- напечатает "[0]" (и даже Warning не вылезет).
        %     * при usefootcite==1, ссылка сработает нормально. В авторской библиографии будут только процитированные в refsection=0 работы.
} % Характеристика работы по структуре во введении и в автореферате не отличается (ГОСТ Р 7.0.11, пункты 5.3.1 и 9.2.1), потому её загружаем из одного и того же внешнего файла, предварительно задав форму выделения некоторым параметрам

\textbf{Объем и структура работы.} Диссертация состоит из~введения,
\formbytotal{totalchapter}{глав}{ы}{}{} и заключения.
Полный объём диссертации составляет
\formbytotal{TotPages}{страниц}{у}{ы}{}, включая
\formbytotal{totalcount@figure}{рисун}{ок}{ка}{ков} и
\formbytotal{totalcount@table}{таблиц}{у}{ы}{}.
Список литературы содержит
\formbytotal{citeexternal}{наименован}{ие}{ия}{ий}.
В главе~\ref{ch:hessian} изучаются спектральные свойства матрицы Гессе нейронных сетей: для полносвязных и сверточных архитектур получены верхние оценки спектральной нормы гаусс"--~ньютоновской компоненты. В главе~\ref{ch:landscapes} исследуется стабилизация ландшафта функции потерь при увеличении объема выборки: на основе локальной квадратичной аппроксимации вводятся критерии стабилизации в одной точке окрестности минимума, по распределению точек в полном пространстве параметров и в подпространстве главных собственных векторов матрицы Гессе. В главе~\ref{ch:scaling-laws} оценки локальной кривизны и критерии стабилизации ландшафта объединяются в критерии достаточного размера выборки и в схему законов масштабирования параметрических моделей. В главе~\ref{ch:linear} вероятностные линейные модели рассматриваются как частный случай параметрических моделей: исследуется стабилизация оценок максимального правдоподобия и апостериорных распределений параметров и строятся критерии достаточного размера выборки. В заключении сформулированы основные результаты работы.